\section{Conclusion}

This study examines OMX Stockholm Benchmark reconstitutions from July 2007 through June 2026 using returns, reported fund holdings and trading data. The main sample contains 256 additions and 215 deletions. The average addition CAR falls from 8.52\% in the first block to 1.50\% in the latest block, but the intermediate periods do not form a monotonic decline. The latest estimate is imprecise, and a failure to reject zero does not establish that event returns have become economically negligible.

The later decade retains a positive addition association of 2.15\%, including under several separately reported exclusions. Benchmark-linked holdings changes and increasingly concentrated trading before effectiveness show that the price pattern occurs alongside measurable investor activity. They do not identify how much of the return is caused by index demand. Early holdings ratios are particularly sensitive to a small and uncertain denominator, while later ratios remain descriptive because coverage and non-index trading are not fully resolved.

An exploratory calculation on fixed, non-overlapping return cohorts finds a positive average addition CAR through twenty trading days after effectiveness in the later decade. The result improves the accounting of the return horizon, but establishes neither a permanent price shift nor causality. Deletions and the latest addition block are less precisely resolved. The defensible conclusion is that Swedish reconstitutions exhibit changing and heterogeneous return associations whose interpretation depends on timing, measurement and inference. A complete disappearance of the effect, a structural demand elasticity and a causal mechanism remain beyond what this design can establish.
